\section{The authority map applied to Mariology}

\begin{longtable}{>{\raggedright\arraybackslash}p{0.19\textwidth}>{\raggedright\arraybackslash}p{0.25\textwidth}>{\raggedright\arraybackslash}p{0.25\textwidth}>{\raggedright\arraybackslash}p{0.23\textwidth}}
\toprule
Class & Marian example & Response & Boundary\\
\midrule
\endhead
Revealed dogma & Divine motherhood, perpetual virginity, Immaculate
Conception, bodily Assumption & Divine and Catholic faith & Exactly four
commonly identified Marian dogmas; their modes and histories differ\\
Authoritative doctrine & Mary's lifelong personal holiness, spiritual
motherhood, subordinate cooperation and intercession & Assent according to the
exact proposal; definitive status may not be invented & Not automatically a
fifth Marian dogma\\
Liturgical and patristic reception & New Eve, All-Holy, Queen, Ark, Dormition
traditions & Reverent reception according to source and ecclesial use &
Typology is not literal history; liturgy does not define every associated
detail\\
Devotional title or practice & Rosary mysteries, approved invocations,
consecration practices & Freedom and prudence within ecclesial approval and
doctrine & Popularity, indulgence, or papal use does not by accumulation define
a dogma\\
Theological opinion & Mary's death, modes of cooperation, interpretations of
symbols & Weight according to evidence and reception & Venerable does not mean
obligatory\\
Private revelation / judgment & Reported apparition, message, shrine, or
pastoral \emph{Nihil obstat} & No divine and Catholic faith; follow competent
discipline and judgment & Does not complete Revelation or authenticate every
message, cure, image, or accretion\\
\bottomrule
\end{longtable}

The companion article \emph{How the Catholic Church Teaches} explains the
universal framework: object, mode, competent intention, and response. This work
does not duplicate that history. It asks what follows when those distinctions
are applied to Marian claims.

